% ============================================================
% Title supplied by appendix wrapper.
\label{sec:formal-certification-boundary}
% ============================================================

This appendix records the role of the Lean formalization relative to the
mathematical claims of the manuscript. It is documentary rather than foundational:
it states what is certified, what is separately audited, and what lies outside
the certified claim.

\subsection{Trusted Meta-Level and Certified Claim}
\label{subsec:formal-trust-boundary}

The trusted meta-level engine is Lean's kernel, dependent type theory,
inductive definitions, definitional equality, and the proof checker itself.
For the system-level architecture of Lean itself, see
\cite{deMouraKongAvigadEtAl2015}.
Inside that engine, the certified claim is intentionally lane-specific: the
active theorem-bearing manuscript surface in Parts~I--II together with the
formal appendix lanes is linked to explicit
audited Lean declarations on designated theorem surfaces, and no hidden
imported mathematical or logical axioms are to be counted as part of those
surfaces.

This is not the claim that every sentence in the repository has been
formalized. The claim is instead that the active formal manuscript items are
tracked against audited Lean declarations on the designated theorem surfaces.

\subsection{Certified Bedrock Lane}
\label{subsec:formal-bedrock-lane}

The primary certified lane is the rebuilt bedrock-and-tower spine. In the
manuscript this is the opening structural chain through Chapters~1--6, now
collected inside Part~I. It starts from the manuscript's lower counting
bedrock and then proceeds forward through the active foundation and
realized-model chain. In the manuscript, that lower bedrock is stated as
Distinction, Free Composition, and Realized Closure, and the older package of
Constancy, Existence, Incidence, and Closure Balance is recovered as the first
theorem ladder. On the current Lean side, however, the certified lane still
begins from the equivalent unpacked declaration surface corresponding to that
recovered package. Acceptance on this lane is based on axiom-surface audits of
the exposed declarations rather than on import inspection alone: the point is
to certify what the exported theorem surface actually depends on.

Accordingly, the bedrock claim is not that Lean contributes extra mathematics.
It is that the active bedrock-and-tower proof chain is forced from the stated
counting bedrock, relative only to Lean's proof checker, with the current Lean
formalization still certifying the equivalent unpacked interface.

\subsection{Separate Lean PartIII Bridge Lane}
\label{subsec:formal-part3-lane}

The continuum-bridge layer at the end of manuscript Part~I is formalized on
a separate Lean root, historically named \texttt{PartIII.lean}, and audited on
its own strict no-axioms surface. That separation is intentional. It allows
the continuum-identification layer to evolve as a bridge from the discrete
calculus into faithful limit classes without altering the bedrock-certified
bedrock-and-tower root.

The Lean PartIII bridge lane is therefore formalized and audited, but it is
not folded back into the narrower statement that the bedrock-and-tower spine
alone is forced from bedrock. The manuscript uses its results as an explicitly
separate bridge lane.

\subsection{Separate Lean PartIV Law-Side Lane}
\label{subsec:formal-part4-lane}

The least-observer-motion layer of manuscript Part~II is also formalized on a
separate Lean root, historically named \texttt{PartIV.lean}, and audited on
its own strict no-axioms surface. That lane sits above the exported
bedrock-and-tower interfaces together with the separate continuum-bridge lane,
together with explicit law-side data for the Part~II physical realization
principle, its observer / exchange / residual-return clauses, the four-step
law-completion chain, and the
non-dictionary sector-generation / classification interfaces used in the
chapter. The current audited surface also includes a concrete finite
observer-motion functional on admissible cuts and a concrete HFT-2-based
returned-flux / unrecovered-stock functional for the residual-return side,
together with a canonical selected Schur bridge whose selected-cut layer
already packages the least-motion cut and residual accounting on one selected
observer datum, and which now reads its law-side content from the grounded
self-sufficient realization bridge on that same datum. On the current Lean
surface, the exact-Schur side is still reduced to the one exact-visible-law
condition that the selected law is already exactly visible on that same
selected cut; the local conditions of observable support and cut commutation
are derived consequences of that clause. The manuscript reads this as the
sharp formal core of the self-sufficient realization principle on the
least-motion cut.

The Lean PartIV law-side lane is therefore covered in Lean as a separate
certified physical-realization/observer/sector-generation/law-completion lane.
But it is not folded back into the narrower statement that the
bedrock-and-tower spine alone is forced from bedrock. The reason for keeping
the Lean PartIV lane separate is the same as for the Lean PartIII lane: the
observer-selection and law-side physical bridge layer should remain explicit rather than being
blurred into the
bedrock-certified root.

\subsection{Separate Appendix Certification Lanes}
\label{subsec:formal-appendix-status}

The supplementary appendices are now covered in Lean as two explicit separate
appendix lanes. The operator-side appendices on advanced operators and
kernel/constraint maps are tracked against a strict no-axioms root
\texttt{AppendixOperators.lean}. The oblique-horizon, continuous spectral
horizon, and trace-level appendices are tracked against a separate strict
no-axioms root \texttt{AppendixHorizons.lean}.

Those appendix lanes are rigorous in the same sense as the other audited
lanes: the theorem-bearing manuscript items are tagged, the tags are checked
against the corresponding Lean roots, and the exposed theorem surfaces are
audited by \verb|#print axioms|. But the appendix lanes are still kept
separate from the narrower bedrock-and-tower bedrock-forcing claim.

The reason for that separation is mathematical honesty. Several appendix
statements require genuinely extra operator, trace, spectral, finite, or
oblique datum. In the rebuilt formalization such ingredients are not imported
silently from archived files or external foundations. They are packaged as
explicit local data on the appendix lanes themselves. Thus the appendices are
covered, but they are not to be described as hidden consequences of the
bedrock spine alone.

\subsection{Manuscript Tags and Audit Discipline}
\label{subsec:formal-tags}

The local \verb|\lean{...}| tags in the manuscript source identify the Lean
declarations intended to certify the corresponding formal manuscript items.
Those tags are hidden in the ordinary reading copy and are intended primarily
for audit builds, coverage scripts, and reproducibility checks.

Their role is a manuscript-to-Lean crosswalk. A tagged item is required to
land on the appropriate audited Lean surface, and the audit reports check both
directions:
\begin{enumerate}[label=(\roman*), itemsep=0.2em]
\item that each active formal manuscript item in Parts~I--II and in the audited appendices carries a nearby Lean tag, and
\item that each tagged Lean name resolves to the correct audited declaration
surface for its lane.
\end{enumerate}
This is a formal-item correspondence discipline, not a literal
sentence-by-sentence certification of every prose clause in the exposition.

In particular, appendix tags now count as certification evidence only when they
resolve to the active audited appendix roots. Historical or exploratory
cross-references elsewhere on disk remain documentary only and are not to be
read as evidence that the corresponding appendix statements are part of the
current certified manuscript surface.

\subsection{Included Lean/mathlib Validation Realizations}
\label{subsec:formal-mathlib-bridge}

The repository also contains an included one-way Lean/mathlib validation lane.
Its role is to show that the exported bridge interfaces can be instantiated
inside standard Hilbert, topological, continuum, and operator settings already
available in mathlib, and thereby to justify the use of the higher-structure
carrier mathematics in the Part~II and appendix interpretations by validating
those standard realizations explicitly. For the architecture and design of the
Lean mathematical library used on that external side, see
\cite{mathlibCommunity2020}. More broadly, this validation lane sits alongside
the recent growth of Lean-backed scientific formalization in adjacent domains,
including chemical physics and high-energy physics
\cite{BobbinEtAl2024,ToobySmith2025}.

That validation lane belongs to the overall formal package and is therefore
shown on the dashboard and in the upload-facing artifacts. But the direction
of use is still one-way: no theorem on the strict manuscript proof lanes
depends on mathlib imports, and those proofs are not discharged by the
existence of the external realizations. The bridge validates external
mathematical realizability; it does not back-feed axioms into the audited
manuscript lanes.

\subsection{Reproducibility}
\label{subsec:formal-reproducibility}

The repository includes build and audit scripts that:
\begin{itemize}[itemsep=0.2em]
\item build the active Lean lanes for the bedrock-and-tower spine, the separate Lean PartIII and PartIV lanes, the two appendix lanes, and the included mathlib validation lane when requested,
\item print axiom dependencies of the exposed theorem surfaces, and
\item regenerate the manuscript tag and coverage reports used to keep the
      manuscript aligned with the formal development.
\end{itemize}

The purpose of those checks is straightforward: as the manuscript and Lean
development change, the theorem-bearing surface, the audited proof surfaces,
and the manuscript tags should remain synchronized. This appendix records the
boundary those checks are meant to protect.
